\documentclass[../../main.tex]{subfiles}
\begin{document}

\begin{flushright}
\footnotesize\textit{【自動文字起こし・要確認】}
\end{flushright}

\noindent
{\bf [解]} 題意から，2平面は $Z = \dfrac{1}{t} Y$, $Z = \dfrac{1}{t} X$ となるから，$Q(t, Y, Z)$ とおける．$RQ$ の傾き及び中点について
\begin{align}
\begin{cases}
\cdot M\Bigl(t, \dfrac{Y+t}{2}, \dfrac{Z}{2}\Bigr) \text{ が } Z = \dfrac{1}{t} Y \text{ 上} \\[1ex]
\cdot \vec{RQ} \cdot \begin{pmatrix} 0 \\ 1 \\ -t \end{pmatrix} = 0
\end{cases}
\dots \begin{cases} Y = \dfrac{t(t^2-1)}{1+t^2} \\[1.5ex] Z = \dfrac{2t^2}{1+t^2} \end{cases} \quad \dots \text{①}
\end{align}
となる．対称性から $R(Y, t, Z)$ となる．この時，$QR$ の中点を $N$ とすると，平面 $OPN$ に関して立体は対称で，$N\Bigl(\dfrac{t+Y}{2}, \dfrac{t+Y}{2}, Z\Bigr)$ となるから，もとめる体積 $V$ として
\begin{align}
V &= 2 \Delta R\text{-}OPN \\
&= 2 \cdot \dfrac{1}{3} \cdot RN \times \Delta OPN \quad (\because OPN \perp RN) \quad \dots \text{②}
\end{align}
と書ける．
\begin{align}
RN &= \sqrt{\Bigl(Y - \dfrac{t+Y}{2}\Bigr)^2 + \Bigl(t - \dfrac{t+Y}{2}\Bigr)^2} \\
&= \dfrac{1}{\sqrt{2}} |Y - t| = \dfrac{1}{\sqrt{2}} \dfrac{2t}{1+t^2}
\end{align}
\begin{align}
\Delta OPN = \dfrac{1}{2} \cdot \sqrt{2}t \cdot Z = \dfrac{1}{\sqrt{2}} \dfrac{2t^3}{1+t^2}
\end{align}
を②に代入して
\begin{align}
V &= \dfrac{2}{3} \cdot \dfrac{1}{\sqrt{2}} \dfrac{2t}{1+t^2} \cdot \dfrac{1}{\sqrt{2}} \dfrac{2t^3}{1+t^2} \\
&= \dfrac{4}{3} \dfrac{t^4}{(1+t^2)^2}
\end{align}

\begin{figure}[htb]
\centering
\begin{tikzpicture}[scale=1.5, >=stealth]
  \draw[->] (0,0,0) -- (2.5,0,0) node[right] {$y$};
  \draw[->] (0,0,0) -- (0,2.5,0) node[above] {$z$};
  \draw[->] (0,0,0) -- (0,0,2.5) node[below left] {$x$};
  \node[below left] at (0,0,0) {$O$};

  \coordinate (P) at (1.5,0,1.5);
  \coordinate (Q) at (1.5,0.8,0);
  \coordinate (R) at (0.8,1.5,0);
  \coordinate (N) at (1.15,1.15,0);

  \draw[dashed] (0,0,0) -- (P);
  \draw[dashed] (0,0,0) -- (N);
  \draw (Q) -- (R);
  \draw (P) -- (Q);
  \draw (P) -- (R);
  \draw (P) -- (N);

  \node[right] at (Q) {$Q$};
  \node[above] at (R) {$R$};
  \node[right] at (N) {$N$};
  \node[below] at (P) {$P$};
\end{tikzpicture}
\caption{四面体 $O$-$PQR$ と中点 $N$ の位置関係}
\end{figure}

\bigskip\noindent{\bf [解2] (後半部)} \\
サラスの公式から
\begin{align}
V &= \dfrac{1}{6} |t Y Z + t Z Y - t Z t - t t Z| \\
&= \dfrac{t}{3} |Y Z - t Z| \\
&= \dfrac{t}{3} \Bigl| \dfrac{t(t^2-1)}{1+t^2} - t \Bigr| \dfrac{2t^2}{1+t^2} \\
&= \dfrac{t}{3} \cdot \dfrac{2t}{1+t^2} \cdot \dfrac{2t^2}{1+t^2} \\
&= \dfrac{4}{3} \dfrac{t^4}{(1+t^2)^2}
\end{align}

\newpage

\end{document}
